\section{Discussion}
\label{sec:discussion}
\emph{Implementation and transfer scope.} Shared seeds and evaluation ceilings do not make constraint handling neutral. Penalty-based methods, box clipping and explicitly constrained SLSQP use different mechanisms and penalty quantities with different units; no normalized-penalty or alternative-repair experiment is supplied. Conclusions therefore concern the method-plus-constraint-handling implementations tested here. The reference DE configuration does not represent the performance of tuned DE or the entire evolutionary-algorithm family. The finer-rose, Gaussian-wake and power-curve studies re-evaluate existing final layouts rather than optimize under the alternative evaluator; they measure transfer of these layouts, not optimizer superiority when search is run under that evaluator. The $5D$/$6D$ and exact-gradient studies are direct re-optimization experiments. In particular, the SSA control conclusion at $4D$/$5D$ reverses at $6D$ and must not be generalized to denser constraints.
In this benchmark, which method appears best depends strongly on the protocol; we discuss why, what follows for hybrid design and layout practice, and condense the protocol into a checklist.

\subsection{What Drives the Conclusions of Metaheuristic Comparisons}
\label{sec:disc-drivers}
\emph{Baseline configuration.} An advantage over PSO should be interpreted against a justified configuration, with its stability assumptions, velocity treatment and boundary handling reported: the old setting lies beyond the order-2 boundary~\eqref{eq:poli}, which marks approximately where this PSO starts to fail, and places PSO at position 5 of Table~\ref{tab:baseline} instead of 1 (Section~\ref{sec:psosetting}). The $c$-sweep links this to Proposition~\ref{prop:pso-stability} without establishing a cause (feasibility starts to decline in the step containing $c^\ast=1.7486$), and velocity clamping largely restores feasibility (92.2\%), whereas zeroing clipped components does not: the setting fails mainly through large steps. The stability region is thus a cheap screen for settings, not a guarantee.
% CHECK-FINAL [C27]: Table baseline: old setting PSO position 5, constriction position 1; CHECK-CSWEEP [S04]/[S05]: feasibility loss begins in the step containing c* and accelerates ("marks approximately where ... starts to fail"; no causal claim, no jump)
% CHECK-THEORY [T01]/[T04]: order-2 violated (4 > 3.50), constriction order-2 stable; CHECK-CSWEEP [S02]: >= \NSFeasBelow % feasible for c <= \NSCBelow; [S03]-[S05]: progressive decline from the step containing c*; [S09]: zeroing clipped velocity no rescue; [S10]/[S11]: vmax clamping \NSFeasVMax % ("largely"), still worse than constriction -> "mainly through large steps"

\emph{Controls.} An advantage of a hybrid's first phase is informative only if it beats random sampling in the farm at equal cost. Square sampling is too weak a control (730 RS-VNS runs draw no feasible sample, mainly because of the spacing), so a salp-swarm phase appears to help against it; against disc sampling only the PSO phase gains at the benchmark spacing (at $6D$ the SSA phase gains as well; Section~\ref{sec:spacing}), and the small SSA gain over RS-VNS is no longer significant with 1$^\circ$ bins (Section~\ref{sec:ablation}). With essentially one compass-search descent as second phase, this measures the value of a start~\cite{Feng2015}.
% CHECK-DIAG [D17] + CHECK-FINAL [C57]/[C60]: square control weak (spacing binding); [C53]/[C54]: no SSA gain over disc sampling; [C55]: LX-SSA-VNS worse; [C56]: PSO-VNS beats RSD-VNS ("only the PSO phase gains"); CHECK-EXTRA [X53]: SSA gain > margin ruled out (seed, case); CHECK-DIR [F17]: SSA-VNS vs RS-VNS n.s. at 1 deg; CHECK-DIAG [D08]: essentially one descent

\emph{Level of equivalence.} Non-significance is not equivalence, and equivalence depends on the level and the evaluation: PSO-VNS and PSO are equivalent at $\pm0.05$~pp across seeds and cases but not under benchmark-group sensitivity, nor when the stored layouts are re-evaluated under a 1$^\circ$ rose, where PSO-VNS has a small significant advantage (Sections~\ref{sec:psovns-pso} and~\ref{sec:robustness}). The levels describe the fixed benchmark and sensitivities to case exchangeability and dependence within its designed groups; they do not validate extrapolation to new farms. The reported minimal margin grows with them (0.026, 0.041 and 0.067~pp); equivalence holds for $N<10$ but not for $N\ge10$, and 19 of 68 cases differ by more than the margin, in both directions.
% CHECK-EXTRA [X39]/[X41]: equivalent at seed and case level, not at cluster level (CR2), minimal margins seed < case < m < cluster; CHECK-DIR [F10]: 1 deg -0.043 pp, 90% CI below 0, not equivalent; CHECK-EXTRA [X44]: N < 10 equivalent, N >= 10 not; [X43]: \NXHetBeyond of \NXHetN beyond +-m, both directions

\emph{Discretization, initialization and budget.} With 1$^\circ$ sub-bins the wake loss of every method rises (PSO-VNS: 1.322\% to 2.335\%), as optimized layouts exploit the gaps between the 15$^\circ$ bins (Section~\ref{sec:robustness}). PSO-VNS and PSO keep the first two places, but MS-SLSQP moves from sixth to third and PSO's lead under the Gaussian wake with 15$^\circ$ bins disappears. From feasible starts, a single local search (VNS alone) ranks first on the six largest cases (Section~\ref{sec:feasinit}), and at 120,030 evaluations the leading methods lie within 0.17~pp (Section~\ref{sec:budget}). A ranking is thus the ranking of a protocol.
% CHECK-DIR [F03]: loss rise of every method (PSO-VNS 1.322 -> 2.335); [F05]: PSO and VNS-based methods rise most; [F06]: PSO-VNS first, PSO second (\NFRankOrderOne); [F07]: MS-SLSQP 6th -> 3rd; [F09]: Gaussian: PSO first at 15 deg only; [F10]: not equivalent at 1 deg
% CHECK-FINAL [C28]: feasible starts, VNS alone first; [C37]/[C45]: 120,030 evaluations, leading methods within \NBudCloseGapOneTwentyK pp of PSO-VNS

\subsection{Implications for the Design of Two-Phase Swarm--Local-Search Hybrids}
\label{sec:disc-hybrids}
\emph{The swarm phase as a feasibility-reaching start.} A first phase need not be a good optimizer, only a cheap source of good starts; in our runs, the PSO phase is valuable largely because it reaches feasibility: at the switch, its global best is feasible in 98.7\% of the runs, against 95.7\% for SSA and 75.7\% for disc sampling, whereas feasible disc samples are not worse (wake loss 1.53\% vs.\ 1.63\%, over different runs). From feasible starts, PSO never improves on the best initial layout and VNS alone ranks first (Section~\ref{sec:feasinit}). A swarm phase is thus most useful when random layouts are infeasible; when a feasible start is cheap (our packing initialization costs no wake-objective evaluations, although its L-BFGS-B packing and geometric checks take time that is not charged to the budget), the budget is better spent on the local search under a model-query budget. The spacing study fits this reading: at $6D$, where every method is feasible less often, even the SSA phase gives a better start than disc sampling (Section~\ref{sec:spacing}), although the experiment does not isolate feasibility as the cause.
% CHECK-FINAL [C16]: PSO switch feasibility higher than SSA / LX-SSA; [C57]: RSD-VNS feasible at the switch more often than RS-VNS; CHECK-DIAG [D18]: PSO > SSA > RS at the switch; switch macros generated (ablation.phase2_loss_reduction_pct); the PSO-VNS vs RSD-VNS switch-loss comparison is descriptive (different run subsets)
% CHECK-FINAL [C28]/[C38]: feasible starts, VNS alone first, PSO and DE never move; CHECK-DIAG [D13]/[D14]: first-move collisions; CHECK-THEORY [T11]

\emph{The local-search phase and the split.} At 6,030 evaluations, the first compass-search descent gives 97.5\% of the Phase-2 gain and is unfinished in 39 of 40 runs with $N\ge12$ (Section~\ref{sec:ablation}). At small budgets the cost of the descent ($2n$ evaluations per step) thus matters more than the neighborhoods, which act only for small $N$; cheaper descents are obvious candidates, but untested. The first descent also repairs infeasible switch points (densest cases: 63.3\% feasible at the switch, all at the end). Of the splits tested, \ensuremath {\omega =0.75} was best and plain PSO ($\omega=1$) had a higher mean loss (Section~\ref{sec:split}): the swarm needs enough evaluations to reach feasibility, and the descent enough to progress.
% CHECK-DIAG [D08]: first descent \NDDescentSharePSOVNS % of the Phase-2 gain; [D09]: N >= 12 unfinished \NDDescentUnfinishedLargePSOVNS of \NDRunsLargePSOVNS; [D10]: accepted cycles only in k = 1; [D11]: infeasible switch points repaired by the first descent; CHECK-FINAL [C24]: dense_500_10 feasible at switch < 100, final 100
% CHECK-FINAL [C47]: 0.75 best (lowest mean loss and rank), omega = 1 higher mean loss than 0.75; [C59]: 0.9 does not beat 0.75 (Section split); CHECK-EXTRA [X36]

\emph{Where PSO-VNS helps and where it does not.} At 6,030 evaluations from random starts, PSO-VNS reaches feasibility more reliably on the Horns Rev~1 block (30/30 vs.\ 19/30 runs, McNemar $p=\ensuremath {9.8\times 10^{-4}}$; on the benchmark PSO is feasible in every run as well) and, added after the primary analysis, on a Lillgrund block at minimum spacing $3D$ (17/30 vs.\ 0/30; Section~\ref{sec:lillgrund}), ranks best on the six largest cases at all three budgets (exploratory), and, in the post hoc subgroup of the larger Data Set~II layouts ($N\ge10$), has the lower mean wake loss in 10 of 10 cases ($p_{\rm Holm}\le0.018$), also with 1$^\circ$ bins. It does not help from feasible starts, in the densest Data Set~I layouts (all 7 Data Set~I cases beyond the margin favor PSO), on IEA37, or on the benchmark average by more than a few hundredths of a percentage point (re-optimized at $5D$ and $6D$, its mean advantage on the split cases is larger: 0.094 and 0.068~pp; Section~\ref{sec:spacing}).
% CHECK-EXTRA [X52]: Horns Rev 30/30 vs 19/30, McNemar p 9.8e-4 (one site, random starts); CHECK-FINAL [C08]/[C20]; generated \NFeasPSO = 100 (R5 O8); [C29]: six largest cases best at all budgets (exploratory)
% CHECK-FINAL [C32] + CHECK-EXTRA [X19]: DS II N >= 10, all-family Holm; CHECK-DIR [F14]: survives 1 deg; CHECK-EXTRA [X43]: all \NXHetDsINPSO DS I cases beyond m favour PSO; [X25]: mean gain < 0.05 % AEP; CHECK-FINAL [C28]/[C40]

\subsection{Implications for Wind Farm Layout Practice}
\label{sec:disc-practice}
The direction discretization is the main caveat for practice: with 1$^\circ$ bins the benchmark wake loss rises by 0.40--1.04~pp, and none of the 901 feasible optimized Horns Rev~1 layouts exceeds the installed block (16 with the 5$^\circ$ bins used in optimization). Gains of layouts optimized under a coarse rose should be verified at a fine resolution before they are credited to an optimizer.
% CHECK-DIR [F03]: every method's loss rises by \NFLossRiseMin to \NFLossRiseMax pp; [F21]: 0 of 901 feasible runs above the installed block at 1 deg (16 at 5 deg); CHECK-FINAL [C62]: no \NHR...Above macro used

The Horns Rev~1 block, without the other 64 turbines, with $4D$ instead of $7D$ spacing and without cables, loads or navigation, tests optimizers, not the installed design (Section~\ref{sec:hornsrev-site}). Our model takes $C_T$ at the free-stream speed (wake loss \ensuremath {-}0.72~pp relative to PyWake's NOJ model); with $C_T$ at the local speed it reproduces PyWake's hub-center NOJ exactly, as design studies require. On IEA37, the best PSO-VNS layouts at 30,030 evaluations lie 1.93\% (16 turbines) and 6.23\% (36 turbines) below the best feasible published layouts in AEP; they came from the participants' own optimizers, some gradient-based, with budgets not reported comparably~\cite{Baker2019,Thomas2023}. On this piecewise-smooth model, exact gradients, alone or after a PSO phase, gave a higher AEP than every metaheuristic at the same charged budget (in objective-equivalent units; per run they took about 3--6 times as long as PSO-VNS) and came within 0.04\% and 1.84\% of the best feasible published layouts (strict; Section~\ref{sec:iea37}); in the comparison of~\cite{Thomas2023}, in which each method was run by researchers experienced with it and without a common budget, gradient-based and gradient-free methods reached similar yields. For smooth, differentiable wake models, our results thus favor gradient-based optimization, possibly started from a swarm, over the metaheuristics compared here; the metaheuristic comparisons of this paper are relevant where the objective is discontinuous, as with the top-hat wake of the benchmark, or where gradients are not available. The methods compared here, at these budgets, are reference points for comparing optimizers, not design tools.
% CHECK-DIR [F32]/[F34]: PyWake 661.39 vs ours 666.75 GWh/yr, wake loss -0.72 pp, cause C_T at free-stream speed; local C_T reproduces PyWake's hub-centre NOJ; [F44]: projected IEA37 gaps -1.93 / -6.23 % (30,030); CHECK-FINAL [C40]: best published layouts better

\subsection{Recommendations for Benchmarking}
\label{sec:disc-recs}
Table~\ref{tab:recommendations} condenses the protocol into a checklist~\cite{BartzBeielstein2020,LaTorre2021}, with the reason for each item in this study and where the paper implements it.
\begin{table}[!t]
\centering
\caption{Checklist for benchmarking metaheuristics on the WFLOP.}
\label{tab:recommendations}
\footnotesize
\begin{tabular}{@{}>{\raggedright\arraybackslash}p{0.02\textwidth}>{\raggedright\arraybackslash}p{0.36\textwidth}>{\raggedright\arraybackslash}p{0.27\textwidth}>{\raggedright\arraybackslash}p{0.25\textwidth}@{}}
\toprule
& Item & Rationale in this study & What we did (where) \\
\midrule
1 & Justified baseline settings, screened with stability conditions where available (PSO: Eq.~\eqref{eq:poli}); state bound and velocity handling & The old setting reverses the order of PSO and the salp-swarm hybrids & $c$-sweep, bound variants (Sections~\ref{sec:pso} and~\ref{sec:psosetting}) \\
2 & Equal model-query budgets (incl.\ initial population, gradients), elapsed time reported separately; seed-paired runs & Iterations differ in cost (LX-SSA: $2N_p$) & $B=6{,}030$; seeds 1--30 (Section~\ref{sec:setup}) \\
3 & Report feasibility rates, conditional quality and an all-run paired outcome separately & Conditional means can favor less reliable methods & Qualification threshold (Table~\ref{tab:friedman68}); all-run endpoint (Section~\ref{sec:archive-reliability}) \\
4 & Paired, multiplicity-corrected tests; robustness to thresholds and clusters & Rank post hoc tests depend on the pool; cases share farms & Holm, six thresholds, six clusters (Section~\ref{sec:setup}) \\
5 & Equivalence at a stated level, margin and evaluator; name the primary questions; report minimal margins & Verdicts change with the level & TOST at three levels (Section~\ref{sec:psovns-pso}) \\
6 & Ablate hybrids against the local method alone and an equal-cost in-site random control & A square control flatters salp-swarm phases & VNS, RS-VNS, RSD-VNS (Table~\ref{tab:ablation}) \\
7 & Vary budget and initialization before generalizing a ranking & The leader changes from feasible starts & Three budgets, feasible starts (Table~\ref{tab:feasbudget}) \\
8 & Re-evaluate at a finer rose and other models, and re-optimize before claiming a robust ranking; publish versioned code and data & Layouts exploit coarse bins & 1$^\circ$ bins, Gaussian, PyWake (Sections~\ref{sec:hornsrev-site} and~\ref{sec:robustness}) \\
\bottomrule
\end{tabular}
\end{table}
% CHECK-DIAG [D01]/[D02]: instrumented runs reproduce the stored runs; CHECK-CSWEEP [S09]/[S10]: bound/velocity handling changes the old setting's feasibility (item 1); CHECK-FINAL [C27] (item 1); CHECK-EXTRA [X02]-[X09], [X11]/[X14], [X18] (item 4); [X41]: minimal margin grows with the level (item 5); CHECK-FINAL [C57]/[C60] (item 6); [C28] (item 7); CHECK-DIR [F03]/[F07] (item 8)

\section{Threats to Validity}
\label{sec:limits}
Each threat below names the claims it limits; unless a check is cited, the threat is disclosed rather than resolved.
\emph{Internal validity.} PSO-VNS repeats PSO exactly up to the switch, and the controls differ from the hybrids only in Phase~1 (3,015 versus 3,030 evaluations); the code reproduces the original and stored runs (Section~\ref{sec:validation}). Configuration remains a threat: DE is untuned (its low feasibility, 71.4\%, may reflect penalty-dominated search), MS-SLSQP uses finite differences on a discontinuous objective (with exact gradients on the piecewise-smooth IEA37 model, SLSQP outperforms every metaheuristic; Section~\ref{sec:iea37}), and the old PSO setting was run without velocity clamping, which largely restores its feasibility, and with one boundary rule and penalty scaling. More generally, the penalty adds boundary violations in m$^2$ to spacing violations in m with $\mu=10^{10}$, and box clipping leaves the circular boundary to the penalty; these choices can interact differently with population moves and coordinate search. No dimensionless-violation or lexicographic feasibility-rule experiment was run, so the weaknesses observed for SSA, LX-SSA, DE and the old PSO setting characterize these implementations under this penalty and repair, not their algorithm families. The SSA--LX-SSA contrast of the main comparison confounds the Laplace step with the doubled cost per iteration; the ablation that separates them (Section~\ref{S-sec:S-rev-lap}) covers only the twelve cases of the split study, on which none of its Holm-adjusted case-mean tests against LX-SSA is significant.
% CHECK (manual, see 05_setup): 2 of 24 Mann-Whitney tests reject; CHECK-DIAG [D01]/[D02]: instrumented runs reproduce the stored runs; CHECK-FINAL [C61]: controls switch at 3,015, swarms at 3,030; generated \NFeasDE ("may reflect"); CHECK-CSWEEP [S10]/[S11]: vmax \NSFeasVMax % ("largely restores"); CHECK-FINAL [C17]: LX-SSA as published

\emph{Construct validity.} At 6,030 evaluations the ``VNS'' phase is essentially one truncated compass-search descent, so the component analysis concerns this local search, not VNS in general. The objective, the gross energy of a discretized rose under a hub-center top-hat Jensen wake, is a benchmark value, not farm quality; final quality at a fixed budget is the only measure tested; and the feasibility-first ranking limits, but does not remove, survivorship bias.
% CHECK-DIAG [D08]-[D10]: truncated compass search; CHECK-FINAL [C21]: cubic curve lowers the power, ranking intact

\emph{External validity.} All results are for discretized roses (15$^\circ$ benchmark bins, 5$^\circ$ on Horns Rev~1), which optimized layouts exploit, a top-hat Jensen wake with the non-physical upstream deficit of the benchmark implementation (Section~\ref{sec:model}), the linear benchmark power curve (Section~\ref{sec:powermodel}), circular farms, a $4D$ spacing (re-optimization at $5D$ and $6D$ covers only the 12 split cases; Section~\ref{sec:spacing}) and a prescribed $N$, with random starts and 6,030 evaluations unless stated: rankings of a small budget (Section~\ref{sec:disc-drivers}). Measured wind climates enter only through two 16-turbine blocks, Horns Rev~1 and, added after the primary analysis, Lillgrund (seven months of measurements, minimum spacing $3D$, 5$^\circ$ bins only; Section~\ref{sec:lillgrund}), on which the feasibility advantage of PSO-VNS rests, with the Horns Rev~1 and IEA37 caveats of Section~\ref{sec:disc-practice}; the 64-turbine IEA37 scenario was not run. The pool lacks CMA-ES~\cite{HansenOstermeier2001}, L-SHADE~\cite{Tanabe2014}, adaptive GAs, ACO~\cite{Eroglu2012} and BBO~\cite{Bansal2017}; DE is the only evolutionary method of the eight-method pool, and a standard real-coded GA, added after the primary analysis (Section~\ref{S-sec:S-rev-ga}), ranks third on the benchmark but best on the Horns Rev~1 block at 30,030 evaluations and, apart from the exact-gradient methods, with 36 IEA37 turbines, so ``best of eight'' means best among the eight methods compared. The exact-gradient baselines were also added after the primary analysis and run only on IEA37 Case Study~1, with one charge per gradient ($c_g=3$; Section~\ref{S-sec:S-rev-grad}). Turbulence and loads (critical at $4D$), wake steering, blockage, terrain, cables, electrical losses and costs are outside our scope; all AEP values are gross.
% CHECK-DIR [F03]-[F07]: coarse bins exploited; [F21]: 0 of 901 feasible runs above the installed block at 1 deg; [F32]/[F34]: PyWake 661.39 vs ours 666.75 GWh/yr, wake loss -0.72 pp, cause C_T at free-stream speed; [F44]: projected IEA37 gaps -1.93 / -6.23 % (30,030)
% CHECK-FINAL [C45] (budget study; NBudCloseGapOneTwentyK generated); CHECK-EXTRA [X52]: Horns Rev feasibility (one site, random starts)

\emph{Statistical conclusion validity.} The margin, the $N\ge10$ split, the Data Set~II subgroup and the density pattern were stated after the primary analysis; the margin is a convention, and the minimal margins let readers apply their own. The 68 cases form only six clusters with nested $N$: an exact sign test over clusters cannot fall below $p=0.031$, cluster-robust intervals have 5 degrees of freedom, and the p-values describe this benchmark, not a population of farms. The budget and feasible-start rankings (six cases) are not tested, and the feasibility result rests on two sites with 30 seeds each, the second added after the primary analysis. The spacing study was added after the primary analysis and rests on 10 split cases (8 for some contrasts at $6D$). Equivalence statements other than the two primary questions (Section~\ref{sec:setup-stats}) are exploratory, and their error rates are not controlled jointly.
% CHECK-FINAL [C49] + CHECK-EXTRA [X41]: post hoc margin, minimal margins; [X47]: density pattern (exploratory); [X10]: six clusters; [X11]: exact p = \NXClMinP; [X42]: CR2 5 d.f.; CHECK-FINAL [C29]: rank differences not tested; CHECK-EXTRA [X52]

\section{Conclusions and Future Work}
\label{sec:conclusion}
On the discretized Kusiak--Song benchmark with random starts and 6,030 evaluations, three findings emerge. First, the baseline configuration changes the conclusions: the old PSO setting, which lies beyond the order-2 boundary and rarely samples feasible layouts late in the run, puts PSO behind salp-swarm hybrids, whereas the constriction setting puts it ahead of every salp-swarm method at this budget. Second, at the benchmark spacing $4D$, salp-swarm first phases give no gain larger than the 0.05~pp margin over random sampling in the farm (a larger gain is ruled out across seeds and cases, not under benchmark-group sensitivity; LX-SSA as published is worse; re-optimized on the split cases, this holds at $5D$ but not at $6D$, Section~\ref{sec:spacing}), whereas a PSO phase gives the second phase, essentially one compass-search descent, a clearly better start. Third, PSO-VNS has the best average rank on the 15$^\circ$ benchmark, and its layouts keep it when re-evaluated with 5$^\circ$ and 1$^\circ$ direction bins, also under a Gaussian wake at 1$^\circ$, but its average advantage over a well-configured PSO is small: within the post hoc margin $\pm0.05$~pp under the benchmark rose (at the seed and case levels, not under scenario-cluster sensitivity), \ensuremath {-}0.043~pp after 1$^\circ$ re-evaluation and \ensuremath {-}0.094 and \ensuremath {-}0.068~pp at minimum spacings $5D$ and $6D$ (split cases), with case-level differences in both directions that follow a data set $\times$ density pattern (exploratory). PSO-VNS reaches feasibility more reliably on Horns Rev~1 (30/30 vs.\ 19/30 runs, McNemar $p=\ensuremath {9.8\times 10^{-4}}$) and on a second measured site, a Lillgrund block (17/30 vs.\ 0/30 and 30/30 vs.\ 19/30 runs at 6,030 and 30,030 evaluations; Section~\ref{sec:lillgrund}); it is a consistent reference method, not a superior optimizer: a real-coded GA, added after the primary analysis, ranks third on the benchmark but does better on the Horns Rev~1 block at 30,030 evaluations and with 36 IEA37 turbines (Section~\ref{S-sec:S-rev-ga}), and on the piecewise-smooth IEA37 model, SLSQP with exact gradients, alone or after a PSO phase, does better than every metaheuristic compared (Section~\ref{sec:iea37}).
% CHECK-FINAL [C12]/[C27]: old setting puts PSO behind salp-swarm hybrids, constriction ahead of every salp-swarm method (6,030); CHECK-DIAG [D05]: few feasible particle evaluations late in the run; CHECK-CSWEEP [S04]
% CHECK-FINAL [C53]/[C54]: no SSA gain over disc sampling = "in the farm"; [C55]: LX-SSA-VNS worse; [C56]/[C14]/[C44]: PSO phase beats both controls; CHECK-EXTRA [X53]: gain > margin ruled out across seeds and cases (one-sided); CHECK-DIAG [D08]: essentially one descent
% CHECK-DIR [F06]/[F09] + CHECK-FINAL [C21]/[C22]: PSO-VNS first at 15/5/1 deg (Jensen), cubic curves, Gaussian 1 deg; PSO first only under Gaussian 15 deg; CHECK-EXTRA [X39]: seed/case equivalent, cluster not; CHECK-DIR [F10]: -0.043 pp at 1 deg; CHECK-EXTRA [X43]: 19/68 beyond margin, both directions; [X47]: data set x density pattern (exploratory); [X52]: Horns Rev 30/30 vs 19/30, McNemar p 9.8e-4; CHECK-FINAL [C08]/[C20]

Future work follows from the threats above. \emph{Expensive wake models}: surrogate-assisted evaluation (graph neural networks mapping a layout and a wind state to turbine powers, or data-driven and Kriging surrogates~\cite{Yang2023,Wang2024,Li2025}) would make equal-budget comparisons at higher fidelity affordable, if the surrogate error is reported next to the method differences. \emph{Larger farms and more sites}, such as the IEA37 64-turbine scenario, the full Horns Rev~1 and Lillgrund farms, and sites with multi-year wind records. \emph{Gradient-based baselines} beyond IEA37 Case Study~1: exact or algorithmically differentiated gradients on smooth models of larger farms~\cite{Rodrigues2024}, and swarm--gradient hybrids, under budgets that charge every gradient. \emph{Stronger baselines and controls}: CMA-ES, L-SHADE, tuned DE, and random-sampling controls at larger budgets. \emph{Finer roses during optimization}, with loads and economic objectives.
